\section{Kết luận và hướng phát triển}
\label{sec:ket-luan}

\subsection{Kết quả đạt được}
Nhóm đã thiết kế và cài đặt hoàn chỉnh một ứng dụng cộng tác (chat chung, nhắn riêng, gửi file) chạy trên giao
thức tự thiết kế trên UDP, chỉ dùng standard library của Go. Những kết quả chính:
\begin{itemize}
  \item \textbf{Tin cậy:} dữ liệu đúng từng byte khi mất 20\% gói mỗi chiều; file 200 MB qua server khớp SHA-256
        mà bộ nhớ server không tăng theo kích thước file.
  \item \textbf{Hiệu năng:} CUBIC dùng hết 2 MB/s của nút cổ chai, quỹ đạo \texttt{cwnd} là đường bậc ba; không
        có HoL: tin chat khi đang truyền file chỉ trễ thêm khoảng 2 ms (p50).
  \item \textbf{Bảo mật:} X25519 + AES-256-GCM, header được xác thực qua AAD, nonce không bao giờ lặp.
  \item \textbf{Ổn định:} 20 client đồng thời với số goroutine cố định; tự kết nối lại; \texttt{CLOSE} khi tắt nên
        danh sách online cập nhật ngay; \fname{go test -race} không phát hiện race condition.
\end{itemize}

\subsection{Hướng phát triển}
Một số hướng có thể mở rộng thêm từ hệ thống hiện tại:
\begin{itemize}
  \item \textbf{Congestion control dựa trên băng thông và RTT} như BBR \cite{bbr} để tránh giảm mạnh \texttt{cwnd}
        khi đường truyền có mất gói ngẫu nhiên (\cref{subsec:cubic-gioi-han}).
  \item \textbf{Xác thực định danh server:} cấu hình khóa ký dài hạn cho server để ký \texttt{server\_pub} trong
        bước \texttt{KEYX}, chống tấn công MITM (\cref{sec:bao-mat}).
  \item \textbf{Lập lịch bằng round-robin giữa các stream file:} tách \texttt{bulkQ} theo từng stream file
        thay vì xếp chung một hàng đợi FIFO, để file nhỏ gửi sau không phải chờ file lớn gửi trước trên cùng kết nối.
  \item \textbf{Truyền file trực tiếp P2P giữa hai client:} dùng server làm kênh báo hiệu để hai client trong cùng
        LAN thực hiện UDP hole punching và gửi thẳng các gói STREAM chứa file cho nhau, giảm tải băng thông chuyển
        tiếp trên server.
\end{itemize}
